% Source-derived chapter 06: The main cohomological results
% Original source: canonical-representations-surface-groups
\chapter{The main cohomological results}
\label{canonical-representations-surface-groups-chapter-06}
\source{canonical-representations-surface-groups}

\begin{remark}[Source section]
\label{canonical-representations-surface-groups-chapter-06-source}
\source{canonical-representations-surface-groups}
\source{canonical-representations-surface-groups}
This chapter reproduces the corresponding section of the retrieved
LaTeX source, including its definitions, statements, examples, and
proofs. It has not yet been translated into Lean.
\notready
\end{remark}

% The source section title was: The main cohomological results
\label{section:vanishing}
\source{canonical-representations-surface-groups}

In \autoref{subsection:rank-bound-proof}, we prove our main cohomological result, showing that higher direct images of 
unitary systems on families of curves contain no low rank local systems.
From this we derive a related vanishing result in \autoref{subsection:no-sections},
which we will use later to establish cohomological rigidity of certain local systems.


Throughout this section, we continue to use \autoref{notation:versal-family}.
Namely, we use $\pi:\mathscr{C}\to \mathscr{M}$ for a versal family of
$n$-pointed curves of genus $g$, and $\pi^\circ: \mathscr{C}^\circ \to \mathscr{M}$ for the
associated punctured versal family of curves.

\subsection{A rank bound}
\label{subsection:rank-bound-proof}
\source{canonical-representations-surface-groups}
We now prove our main result on the cohomology of unitary local systems on versal families of curves,
\autoref{theorem:rank-bound-subsystem}.
That is, for $\mathbb V$ a unitary local system on $\mathscr C^\circ$, we
will show that any nonzero
sub-local system of $R^1 \pi^\circ_* \mathbb V$ has rank at
least $2g-2\rk\mathbb{V}$.

\begin{proof}[Proof of \autoref{theorem:rank-bound-subsystem} assuming
	\autoref{lemma:low-rank-vector} (below)]
Let $m \in \mathscr{M}$, $C=\pi^{-1}(m)$, $C^\circ = (\pi^{\circ})^{-1}(m)$
and $D := C - C^\circ$.
Let $(E, \nabla)$ denote the
Deligne canonical extension
to $C$ of the vector bundle $\mathbb V|_{C^\circ} \otimes \mathscr O_{C^\circ}$ with its tautological connection; $(E,\nabla)$ is a flat vector bundle on $C$ with regular singularities along $D$. The residues of $\nabla$ endow $E$ with the structure of a parabolic bundle $E_\star$, as described in \cite[Definition 3.3.1]{LL:geometric-local-systems}.
By the Mehta-Seshadri correspondence \cite{mehta1980moduli},
the parabolic bundle $E_\star$ is a sum of parabolically stable bundles, hence
parabolically semistable. 

The connection $\nabla$ on $\mathscr{H} := R^1 \pi^\circ_*\mathbb V \otimes_{\mathbb C}
\mathscr O_{\mathscr M}$ induces an $\mathscr O_{\mathscr M}$-linear map
\begin{align*}
	F^1\mathscr{H} \xrightarrow{\overline{\nabla}} \mathscr{H}/F^1 \mathscr{H} \otimes \Omega_{\mathscr M}.
\end{align*}
as in \eqref{equation:shift-derivative}.
\autoref{theorem:period-map-computation} identifies the fiber of this map at the point $m \in \mathscr M$ with the map
\begin{align*}
	\mu: H^0(C, E \otimes \omega_C(D)) &\to \Hom(H^0(C, E^\vee \otimes \omega_C),
	H^0(C, \omega_C^{\otimes	2}(D))) \\
v&\mapsto (u\mapsto B_E(v, u)) \end{align*}
defined via $B_E$ from \eqref{equation:bilinear-pairing}. For $v\in H^0(C, E\otimes \omega_C(D))$, we denote by $\overline{\nabla}_m(v)$ the induced map 
\begin{align*}
 \overline{\nabla}_m(v): H^0(C, E^\vee \otimes \omega_C)& \rightarrow H^0(C, \omega_C^{\otimes	2}(D))\\
u & \mapsto B_E(v, u).
\end{align*}
%$$\overline{\nabla}_m(v): H^0(C, E^\vee \otimes \omega_C)\to H^0(C,
%\omega_C^{\otimes	2}(D))$$ $$u\%mapsto B_E(v, u).$$

Our strategy is to observe that a sub-local system of $R^1\pi_*^\circ \mathbb{V}$ of low rank would provide us with a vector $v$ in $F^1\mathscr{H}_m$ so that the linear transformation $\overline{\nabla}_m(v)$ has low rank. This is in tension with \autoref{proposition:pairing-result} above, which will provide us with the required dimension bound.

Now, suppose we have a nonzero irreducible sub-local system $\mathbb L \subset R^1 \pi^\circ_* \mathbb
V$.
We will show in \autoref{lemma:low-rank-vector} that, after possibly replacing $\mathbb{V}$ with its complex conjugate $\overline{\mathbb{V}}$, there is some nonzero $v \in
F^1\mathscr{H}_m \simeq
H^0(C, E \otimes \omega_C(D))$ for which
$\overline{\nabla}_m(v)$ has rank (as a linear map) at most $\frac{\rk \mathbb L}{2}$. 
Under the identification of
\autoref{theorem:period-map-computation},
we obtain that $\mu(v)$ has rank at most $\frac{\rk \mathbb L}{2}$. 
Using \autoref{proposition:pairing-result}, $\rk \mathbb{V}=\rk
E \geq g-\frac{\rk \mathbb L}{2}$. Upon rearranging, we find $\rk\mathbb{L}\geq 2g-2\rk\mathbb{V}$.
\end{proof}

We now complete the proof of \autoref{theorem:rank-bound-subsystem} by verifying
the following lemma, which is essentially an application of the theorem of the
fixed part.
\begin{lemma}
\source{canonical-representations-surface-groups}
\notready\label{lemma:low-rank-vector}
\source{canonical-representations-surface-groups}
Suppose $\pi^\circ:\mathscr{C}^\circ\to \mathscr{M}$ is a punctured versal family of genus
$g$ curves and $\mathbb{V}$ is a unitary local system on $\mathscr{C}^\circ$. 
For $m \in \mathscr M$, let $\overline{\nabla}_m$ denote the fiber over $m$ of
the map $F^1\mathscr{H} \xrightarrow{\overline{\nabla}} \mathscr{H}/F^1
\mathscr{H} \otimes \Omega_{\mathscr M}$
from \eqref{equation:shift-derivative}.
Suppose 
$\mathbb L \subset R^1 \pi^\circ_* \mathbb V$
is a nonzero irreducible sub-local system.
After possibly replacing $\mathbb{V}$ and $\mathbb L$ with their complex
conjugates
$\overline{\mathbb{V}}$ and $\overline{\mathbb L}$, there is some nonzero $v \in \mathbb L_m \cap F^1
\mathscr H_m$ so that
$\overline{\nabla}_m(v)$ has rank at most $\frac{\rk \mathbb L}{2}$.
\end{lemma}
\begin{proof}
For a local system $\mathbb W$, let $\widetilde{\mathbb W}$ denote the
corresponding local system with real structure as in
\eqref{equation:real-construction}. Note that $$\widetilde{R^1 \pi^\circ_* \mathbb{V}}=R^1 \pi^\circ_* \widetilde{\mathbb{V}}.$$

By \autoref{proposition:subrep-real-variation}, $\widetilde{\mathbb{L}}$
underlies a real polarizable variation of Hodge structure, and we have a
nonzero real mixed Hodge structure $Q_{\mathbb R}$ and a nonzero map of real
variations of mixed Hodge structures
$$\iota: Q_{\mathbb R} \otimes \widetilde{\mathbb L} \to \widetilde{R^1 \pi^\circ_* \mathbb V}.$$  
%%Here is an alternate conclusion of the proof:
%Using the decomposition of \autoref{lemma:three-part-hodge-structure}, any such
%map must also be nonzero on the first part of the Hodge filtration, so we obtain a nonzero map 
%$$F^1(Q \otimes \widetilde{\mathbb L})_{\mathbb{C}} \to F^1(\widetilde{R^1 \pi^\circ_* \mathbb
%V})_{\mathbb{C}}.$$
%Since $Q_\mathbb{C}=\oplus Q^{i,j}$ is (the complexification of) a real mixed Hodge structure itself, we can find some $i,j$ and some $1$-dimensional
%homogeneous subspace $P = \on{Span}(p) \subset Q^{i,j}$, for which 
%$F^1(P \otimes \widetilde{\mathbb L})_\mathbb{C} \to F^1(\widetilde{R^1 \pi^\circ_* \mathbb
%V})_\mathbb{C}$
%is nonzero. Note that $P\otimes \mathbb{L}$ underlies a complex PVHS. 
%Further, as $\mathbb{L}$ is irreducible, the natural map $$P\otimes \mathbb{L}\to (\widetilde{R^1 \pi^\circ_* \mathbb
%V})_\mathbb{C}$$ is an injective map of complex PVHS's. Moreover note that $P\otimes \mathbb{L}$ is supported in at most two bidegrees, as $\mathbb{L}$ appears in $R^1 \pi^\circ_* \mathbb V$.
%
%After possibly replacing $\mathbb L$ with $\overline{\mathbb L}$ and $\mathbb{V}$ with $\overline{\mathbb{V}}$, 
%we may assume both
%$$\dim F^1(P\otimes \mathbb{L})\geq \dim P\otimes
%\mathbb{L}/F^1(P\otimes \mathbb{L}),$$ and,
%using \autoref{lemma:three-part-hodge-structure}, that
%there is an injective map
%\begin{align*}
%\iota: F^1(P \otimes \mathbb L) \to F^1(R^1 \pi^\circ_* \mathbb V).
%\end{align*}
%Set $\mathbb N := P \otimes \mathbb L$.
%Recall $m \in \mathscr M$ is our chosen basepoint.
%Taking any $\ell \in F^1\mathbb N_m$ we set $v = \iota_m(p \otimes \ell)$, for
%$P = \on{span}(p)$ as above.
%
%To conclude, we verify $\rk \overline\nabla_m(v) \leq \frac{\rk \mathbb L}{2}$.
%Since $\iota_m$ is injective,
%it is enough to show
%$\rk \overline \nabla_{\mathbb N, m}(p \otimes \ell) \leq \frac{\rk \mathbb L}{2}$,
%for
%\begin{align*}
%	\overline \nabla_{\mathbb N, m}: F^1\mathbb N_m \to \left( \mathbb N_m /F^1 \mathbb N_m\right) \otimes \Omega^1_{\mathscr{M},m}
%=\Hom((\mathbb N_m / F^1 \mathbb N_m)^\vee, \Omega^1_{\mathscr{M},m})
%\end{align*}
%the derivative of the period map associated to $\mathbb N$ at $m$.
%This holds because $\rk \mathbb N_m/ F^1 \mathbb N_m \leq \frac{\rk \mathbb L}{2}$,
%since 
%$\rk F^1 \mathbb N \geq \rk \mathbb N/F^1 \mathbb N$
%and  $\rk F^1 \mathbb N + \rk \mathbb N/F^1 \mathbb N = \rk \mathbb N = \rk
%\mathbb L.$
%Fix $p \in Q$ for which $\iota(p \otimes \mathbb L)$ is nonzero
%and let $P := \on{Span}(p)$ and let $\mathbb N := P \otimes \mathbb L$.
%We may assume $\rk F^1(\mathbb N) \geq \mathbb N/F^1 \mathbb N$,
%as it holds after possibly replacing $\mathbb L$ and $\mathbb V$ with their
%conjugates by \eqref{equation:dual-conjugate}.
%By irreducibility of $\mathbb N$, $\iota_m$ is
%injective on $\mathbb N$.
%Choose any nonzero $\ell \in \mathbb L_m$ and take $v := \iota_m(p \otimes
%\ell)$.
%In order to show $\mu(v) = \mu(\iota_m(p \otimes \ell))$ has rank at most $\frac{\rk \mathbb
%L}{2}$,
%it is enough to show $\rk \mu'(p \otimes \ell) \geq \frac{\rk \mathbb L}{2}$,
%for
%\begin{align*}
%	\mu': F^1\mathbb N_m \to \left( \mathbb N_m /F^1 \mathbb N_m\right) \otimes \Omega^1_{\mathscr{M},m}
%=\Hom((\mathbb N_m / F^1 \mathbb N_m)^\vee, \Omega^1_{\mathscr{M},m})
%\end{align*}
%the period map associated to $\mathbb N$.
%This holds because $\rk \mathbb N_m/ F^1 \mathbb N_m \leq \frac{\rk \mathbb L}{2}$,
%by assumption that 
%$\rk F^1 \mathbb N \geq \rk \mathbb N/F^1 \mathbb N$
%and  $\rk F^1 \mathbb N + \rk \mathbb N/F^1 \mathbb N = \rk \mathbb N = \rk
%\mathbb L.$
%We now produce the desired $v$. By \autoref{corollary:conjugation-trick},
%$\mathbb{L}$ underlies a complex variation of Hodge structure, and there exists
%a bigraded vector space $Q$ and a nonzero flat map $$\iota: Q\otimes
%\mathbb{L}\otimes \mathscr{O}_\mathscr{M}\to
%\mathscr{H},$$ respecting the bigradings on both sides. 

Let $Q := (Q_\mathbb R) \otimes_{\mathbb R} \mathbb C$.
As the grading on $(\widetilde{R^1 \pi^\circ_* \mathbb V})_{\mathbb C}$ is supported in degrees $(1,0), (0,1), (1,1)$, we may assume (after regrading and possibly replacing $Q$ with a subspace) that either
\begin{enumerate}
\item the bigrading on $\mathbb{L}$ is supported in degree $(0,0)$ and the bigrading on $Q$ has support contained in  $\{(1,0), (0,1), (1,1)\}$, or
\item the bigrading on $\mathbb{L}$ has support equal to $\{(1,0),(0,1)\}$
	and the bigrading on $Q$ is supported in degree $(0,0)$.
\end{enumerate}
In the two cases above, we choose $v$ as follows.

\emph{Case (1):} After possibly replacing ${\mathbb{V}}, \mathbb{L}$ with their complex conjugates, we may assume that  $$\iota_{i,j}: Q^{i,j}\otimes\mathbb{L}\to R^1\pi_*^\circ \mathbb{V}$$ is nonzero for some $(i,j)\in \{(1,0), (1,1)\}$. Choose $\ell\in \mathbb{L}_m$ and $q\in Q^{i,j}$, $(i,j)\in\{(1,0), (1,1)\}$ such that $v:=\iota_m(q\otimes\ell)$ is nonzero.

\emph{Case (2):} After possibly replacing $\mathbb{V}$ and $\mathbb{L}$ with
their complex conjugates, we may assume $\dim\mathbb{L}_m^{0,1}\leq \dim \mathbb{L}_m^{1,0}$. 
Choose arbitrary nonzero $q\in Q$ and $\ell\in \mathbb{L}_m^{1,0}$ such that $v:=\iota_m(q\otimes\ell)$ is nonzero.

It is enough to argue that for $v$ as above, $\overline{\nabla}(v)$ has rank at
most $\frac{\rk\mathbb{L}}{2}$. Since $\mathbb{L}$ carries a complex PVHS, we
may compute the rank of $\overline{\nabla}(v)$ by computing with the analogous
period map 
\begin{align*}
\overline{\nabla}_{\mathbb L,m}: F^1\mathbb{L}_m\to \mathbb{L}_m/F^1\mathbb{L}_m\otimes \Omega^1_{\mathscr{M},m}
\end{align*}
associated to $\mathbb{L}$. In case (1), this map is identically
zero, so we are done (i.e.~$\rk \overline{\nabla}(v)=0$). In case (2), 
from the definition of $\overline{\nabla}_{\mathbb L,m}$,
it
is enough to show
\begin{align*}
	\dim \mathbb{L}_m/F^1\mathbb{L}_m\leq \frac{\rk \mathbb L}{2}.
\end{align*}
This follows because we have 
$\dim \mathbb{L}_m=\dim F^1 \mathbb{L}_m+ \dim \mathbb{L}_m/F^1 \mathbb L_m$ and
we assumed
\begin{equation*}
	\dim F^1 \mathbb L_m = \dim \mathbb L^{1,0}_m \geq \dim \mathbb
	L^{0,1}_m = \dim \mathbb L_m/F^1 \mathbb L_m.\qedhere
\end{equation*}
\end{proof}

\subsection{A vanishing result}
\label{subsection:no-sections}
\source{canonical-representations-surface-groups}

We apply \autoref{theorem:rank-bound-subsystem} to prove a vanishing result about local systems on $\mathscr{C}^\circ$ that are not necessarily unitary. We assume only that their restriction to each fiber of $\pi^\circ$ is unitary. The reader may find it useful to consider the case where $A=\mathbb{C}$ below, though we will need the full generality of the result below for our deformation-theoretic applications.
\begin{theorem}
\source{canonical-representations-surface-groups}
\notready
	\label{theorem:unitary-rigid}
\source{canonical-representations-surface-groups}
	With notation as in \autoref{notation:versal-family}, 
	%let $\pi: \mathscr{C}\to \mathscr{M}$ be a versal family of $n$-pointed curves of genus $g$, with associated family of punctured curves 
	let $\pi^\circ: \mathscr{C}^\circ\to \mathscr{M}$ be a punctured versal
	family of genus $g$ curves. Let $m\in \mathscr{M}$ be a point and set $C^\circ
	= \mathscr C^\circ_m$. Let $A$ be a local Artin $\mathbb{C}$-algebra with residue field $\mathbb{C}$, and let $\mathbb V$ be a locally constant sheaf of free $A$-modules on $\mathscr C^\circ$
	such that
	\begin{enumerate}
 \item $\mathbb V|_{C^\circ}$ is a constant deformation of a unitary local system, i.e. there exists a unitary complex local system $\mathbb{V}_0$ on $C^\circ$ and an isomorphism $\mathbb V|_{C^\circ}\simeq \mathbb{V}_0\otimes A$, and
\item $\rk_A \mathbb V < g$.
		\end{enumerate}
 Then
	$H^0(\mathscr M, R^1 \pi^\circ_* \mathbb V) = 0$.
\end{theorem}
\begin{proof}
	Given a dominant map $\mathscr M' \to \mathscr M$,
	we have an injection $$H^0(\mathscr M, R^1 \pi^\circ_* \mathbb
	V) \to H^0(\mathscr M', R^1 \pi^\circ_* \mathbb V|_{\mathscr M'}).$$
	Hence, it is enough to prove the claim after replacing $\mathscr M$ by
	some $\mathscr M'$ with a dominant map to  $\mathscr M$.
	Combining this with \autoref{lemma:unitary-direct-sum}, we may assume $\mathbb V \simeq \oplus_{i=1}^s \mathbb U_i
\otimes (\pi^\circ)^* \mathbb
	W_i$, where $\mathbb W_i$ are locally constant sheaves of free $A$-modules on $\mathscr M$ and
	$\mathbb U_i$ are unitary complex local systems on $\mathscr C^\circ$. 
	It is enough to show that 
	$$H^0(\mathscr{M}, R^1\pi^\circ_*(\mathbb{U}_i\otimes
	(\pi^\circ)^*\mathbb{W}_i))=H^0(\mathscr{M},
	(R^1\pi^\circ_*\mathbb{U}_i)
	\otimes \mathbb{W}_i )=0$$ for each $i$. 
	Letting $\mathfrak{m}_A$ be the maximal ideal of $A$, the $\mathfrak{m}_A$-adic filtration on $\mathbb{W}_i$ has associated-graded pieces isomorphic to direct sums of copies of $\mathbb{W}_i^0:=\mathbb{W}_i\otimes_A\mathbb{C}.$ 
	Hence $\mathbb W_i\otimes R^1\pi_*\mathbb{U}_i$ has a filtration with associated-graded pieces isomorphic to $\mathbb{W}_i^0\otimes R^1\pi_*\mathbb{U}_i$. 
	Thus it suffices to prove  $$H^0(\mathscr{M},
	R^1\pi^\circ_*(\mathbb{U}_i\otimes
	(\pi^\circ)^*\mathbb{W}^0_i))=H^0(\mathscr{M}, (R^1\pi^\circ_*\mathbb{U}_i) \otimes \mathbb W^0_i)=0.$$
	 Note that $\on{rk} \mathbb{U}_i\otimes (\pi^\circ)^*\mathbb{W}^0_i<g$.

	It is therefore enough to show that if $\mathbb{U}$ is a unitary complex local system on
	$\mathscr C^\circ$ and $\mathbb{W}$ is an arbitrary complex local system on $\mathscr M$
	with
\begin{align*}
	H^0(\mathscr{M}, R^1\pi^\circ_* (\mathbb{U}\otimes (\pi^\circ)^*\mathbb{W}))=
	H^0(\mathscr M, R^1 \pi^\circ_* \mathbb U \otimes \mathbb W) \neq 0,
\end{align*}
	then $\rk \mathbb U \otimes (\pi^\circ)^*\mathbb W \geq g$.
A nonzero element of $H^0(\mathscr M, R^1 \pi^\circ_* \mathbb U \otimes \mathbb W)$ yields a nonzero map of local systems $$\mathbb{W}^\vee\to R^1\pi^\circ_*\mathbb{U}.$$
Since $\mathbb U$ is unitary, we have
$\rk \mathbb{W} \geq 2g - 2\rk \mathbb U$
by \autoref{theorem:rank-bound-subsystem}. 
Hence $\rk \mathbb{W} + 2\rk \mathbb{U} \geq 2g$.
As $\rk \mathbb{W}, \rk\mathbb{U}$ are positive integers, we find 
$\rk \mathbb W \cdot \rk \mathbb U$ is minimized when $\rk \mathbb W = 1$
in which case $\rk \mathbb W \cdot \rk \mathbb U \geq \lceil (2g-1)/2 \rceil \geq g$. 
Note here we use that $g \geq 2$ in order for condition $(2)$ of the statement to be
satisfied as $\mathbb V \neq 0$.
%To see this, we want to minimize the value of $\rk \mathbb W \cdot 2 \cdot(g -
%\rk \mathbb W)$.
%This corresponds to a downward sloping parabola, and the minimum therefore
%occurs at one of the endpoints, which are $\rk \mathbb W = 1$ and $\rk \mathbb
%W = g$. The second case attains the minimum when $g > 1$, and the first case is
%not possible when $g \leq 1$.
In general, this implies
\begin{align}
	\label{equation:rank-product-bound}
\source{canonical-representations-surface-groups}
\rk (\mathbb{U}\otimes (\pi^\circ)^*\mathbb{W})=\rk \mathbb W \cdot \rk \mathbb U \geq g,
\end{align}
as desired.
\end{proof}
%\begin{remark}
%	Note that equality in \eqref{equation:rank-product-bound} occurs only
%	when $\rk \mathbb U = 1$.	
%	If we knew for some reason that the relevant $\mathbb U$ satisfied
%	$\rk U > 1$ we could obtain the stronger bound $\rk \mathbb V < 2g-2$.
%\end{remark}
